\section{Linear Block Codes and Error-Correcting Codes}
\label{sec:bg-ecc}

Error control coding is a standard technique for improving the
reliability of memory systems. In OxRAM memories, the raw bit error rate
(RBER) is lower than in conventional 3D-NAND flash and is comparable to
that of DRAM, making simple linear block codes generally sufficient to
achieve an acceptable post-correction BER in the range of
$10^{-8}$--$10^{-12}$~\cite{patel2020bit}. Several studies have shown
that the optimal error control strategy for OxRAM balances ECC strength
against WRITE-VERIFY complexity, since stronger ECCs permit looser
programming tolerances while simpler WRITE-VERIFY schemes reduce
programming time~\cite{hayakawa2017resolving,
	fukuyama2019comprehensive}. While the exact ECC implementations used in
commercial products are considered vendor trade
secrets~\cite{patel2020bit}, available evidence from manufacturers
suggests that simple linear block codes are sufficient for commercial
OxRAM devices~\cite{hayakawa2017resolving}.

A linear block code (LBC) protects a $k$-bit dataword by appending $r$
parity bits to form an $n$-bit codeword, where $n = k +
	r$~\cite{hamming1950error, sklar2004abcs}. Encoding is performed by
multiplying the dataword $\mathbf{d} \in \mathrm{GF}(2)^k$ by a
generator matrix $\mathbf{G} \in \mathrm{GF}(2)^{k \times n}$:

\begin{equation}
	\label{eqn:codeword-gen}
	\mathbf{c} = \mathbf{d}\mathbf{G}.
\end{equation}

In the canonical systematic form, the generator matrix is structured as
$\mathbf{G} = [\mathbf{I}_k \mid \mathbf{P}]$, where $\mathbf{I}_k$ is
the $k \times k$ identity matrix and $\mathbf{P}$ is the $k \times r$
parity submatrix. In systematic form, the codeword is simply the data
word concatenated with the parity bits:

\begin{equation}
	\mathbf{c} = \mathbf{d}[\mathbf{I}_k \mid \mathbf{P}] = [\mathbf{d} \mid \mathbf{d}\mathbf{P}].
\end{equation}

The parity submatrix $\mathbf{P}$ encodes the linear relationships
between data bits and parity bits: each column of $\mathbf{P}$ specifies
which data bits are XORed to produce the corresponding parity bit. The
code rate $R = k/n$ expresses the fraction of the codeword that carries
user data.

Decoding is performed using the parity check matrix $\mathbf{H} =
	[\mathbf{I}_{n-k} \mid \mathbf{P}^T]$, which is orthogonal to
$\mathbf{G}$: $\mathbf{G}\mathbf{H}^T = \mathbf{0}$. Given a potentially
corrupted codeword $\mathbf{r}$ read from memory, the syndrome is
computed as:

\begin{equation}
	\label{eqn:codeword-syndrome}
	\mathbf{S} = \mathbf{r}\mathbf{H}^T.
\end{equation}

If no errors occurred, $\mathbf{S} = \mathbf{0}$. A non-zero syndrome
indicates that one or more errors are present. The structure of
$\mathbf{H}$ determines whether detected errors can be corrected: a code
with minimum Hamming distance $d_\text{min}$ between codewords can
correct up to $t = \lfloor(d_\text{min} - 1)/2\rfloor$ bit errors per
codeword~\cite{hamming1950error, sklar2004abcs}.

Single-error-correction, double-error-detection (SEC-DED) codes are
commonly used in DRAM and embedded memory systems. Correcting a single
bit error requires encoding the position of the error, which demands at
least $\lceil\log_2(n)\rceil$ parity bits; the double-error detection
capability adds one additional parity bit~\cite{sklar2004abcs}. More
powerful codes such as BCH codes offer higher correction capability at
the cost of additional parity overhead and more complex encoding and
decoding circuitry. The appropriate code strength for a given memory is
determined by the target post-correction BER and the RBER of the
underlying array~\cite{hayakawa2017resolving,
	fukuyama2019comprehensive}.

In commercial ReRAM products, the parity submatrix $\mathbf{P}$ is not
disclosed by the vendor. Without knowledge of $\mathbf{P}$, the parity
bit transitions associated with any given data write cannot be
predicted, and codeword-level timing models cannot be constructed. The
recovery of $\mathbf{P}$ from write-access timing is the central
contribution of Chapter~\ref{chap:brewrc}.